\documentclass[10pt,a4paper]{article}
\usepackage[margin=23mm]{geometry}
\usepackage{amsmath,amssymb,booktabs,tabularx,array}
\usepackage{xcolor,listings}
\usepackage[hidelinks]{hyperref}
\definecolor{codebg}{RGB}{247,248,250}
\lstset{basicstyle=\ttfamily\fontsize{8.5}{9.5}\selectfont,
backgroundcolor=\color{codebg},frame=single,rulecolor=\color{black!20},
breaklines=true,columns=fullflexible,keepspaces=true,
showstringspaces=false,comment=[l]{\#},commentstyle=\color{black!55}}
\newcommand{\ind}{\mathbf 1}
\newcommand{\clip}{\operatorname{clip}}
\setlength{\parindent}{0pt}
\setlength{\parskip}{5pt}
\setlength{\emergencystretch}{2em}
\title{PPO for Tower Selection and Handover Reservations\\
\large High-Level System Design, Pseudocode, and Reward}
\author{ProRAN Project}
\date{8 October 2026}
\begin{document}
\maketitle
\setlength{\abovedisplayskip}{5pt}
\setlength{\belowdisplayskip}{5pt}
\setlength{\abovedisplayshortskip}{3pt}
\setlength{\belowdisplayshortskip}{3pt}

\section*{1. What the system does}
The controller chooses a suitable tower for each device and reserves capacity before a likely handover. It considers signal quality, free resources, load, estimated latency, and the device's service requirements. A strong signal alone does not make a tower the best choice.

\textbf{Inputs:} device position and velocity, current tower, demand, service type, tower capacity and load, active reservations, and estimated latency.

\textbf{Actions:} stay or switch tower; reserve or do nothing; hold half or full predicted demand; activate now, halfway to arrival, or just before arrival.

\section*{2. What model is trained?}
\textbf{PPO is the reinforcement-learning algorithm that trains the actors and critic.} The system design uses two linear actor heads and one linear critic:

\renewcommand{\arraystretch}{1.15}
\begin{tabularx}{\linewidth}{@{}>{\raggedright\arraybackslash}p{0.25\linewidth}X@{}}
\toprule Model & Learned role \\ \midrule
Connection actor & Choose an eligible tower or stay connected to the current one. \\
Reservation actor & Choose the future reservation tower, amount, and activation time. \\
Critic & Estimate future reward from the network state, helping the actors learn. \\
\bottomrule
\end{tabularx}

A linear actor computes action scores from features and learned weights:
\begin{equation}
\mathbf z=W\mathbf o+\mathbf b,
\qquad
\pi(a\mid\mathbf o)=\frac{e^{z_a}}{\sum_{j\in\text{legal actions}}e^{z_j}}.
\end{equation}
$\mathbf o$ is the observation, $W,\mathbf b$ are learned parameters, and $\pi$ is the action probability. Illegal actions have probability zero. Each actor has its own weights and action mask. There are no hidden neural-network layers in this design.

The critic predicts a value from global network state $\mathbf S$:
\begin{equation}
V(\mathbf S)=\mathbf w_V^T\mathbf S+b_V.
\end{equation}
Its weights are trained to predict the reward obtained over time. It does not directly choose a tower.

\textbf{Training:} run simulated episodes, sample actions, observe service and waste, calculate rewards, and update weights using PPO. Optional rule imitation can initialize the actors. \textbf{Normal operation:} load the saved weights and choose the highest-probability legal actions, without updating weights.

\textbf{Prediction is separate:} project movement using current velocity. This is not a trained predictor and can be wrong when movement or traffic changes.

\newpage
\section*{3. One high-level pseudocode}
At most one new reservation is requested per interval. Connection decisions are processed by service priority and device ID, using updated capacity after earlier decisions.
\begin{lstlisting}
IF training: initialize both actors and critic
ELSE: load saved actor weights
FOR each episode or normal run:
    Reset devices, traffic, allocations, and reservations
    REPEAT every interval:
        Store the observable global network state
        waste = 0; successful_switches = 0
        # A. Prepare resources for a likely future handover
        Project movement and demand over the forecast horizon
        Select an uncontracted device likely to need a handover
        IF a candidate exists:
            Observe ALL candidate towers, not just the closest
            Mask invalid holds; keep NO RESERVATION legal
            Reservation actor chooses tower, amount, and timing
            Create a pending contract if a hold is chosen
        FOR each pending contract due now:
            IF its whole amount fits current free capacity:
                Activate it and hold the resources
            ELSE: reject it; consume no capacity
        # B. Serve existing devices, then move them
        Allocate ordinary service using capacity minus active holds
        Record held capacity for the interval's holding penalty
        Move devices with mobility noise and boundary reflection
        # C. Choose actual connections after observing movement
        FOR each device in priority / ID order:
            Refresh tower load, free capacity, signal, and latency
            Mask destinations with poor coverage or too few RBs
            Keep STAY legal; discourage rapid repeated switching
            Connection actor chooses current tower or destination
            IF destination differs from current tower:
                Recheck coverage, capacity, and radio outcome
                available = matching ACTIVE hold + free capacity
                possible_grant = min(actual demand, available)
                admitted = covered AND minimum RBs met AND radio_ok
                IF admitted:
                    Release old allocation; settle device contracts
                    Grant resources and commit the new association
                    successful_switches += 1
                ELSE:
                    Keep old association and allocation
                    Settle attempted contracts; reserved RBs used = 0
                Count unused active reservation RBs as waste
                Check full demand AND latency for handover QoS
        # D. Measure service and reward
        Expire remaining due contracts AFTER connection attempts
        Add their unused active RBs to waste
        For every device, effective grant = actual grant if covered
                                            else zero
        Compute service F, QoS Q, unmet demand U, and under-service O
        Compute switching H, waste W, and held-capacity R penalties
        reward = wf*F + wq*Q - wu*U - wo*O - wh*H - ww*W - wr*R
        IF training: store actions, masks, log probabilities, reward
        Assert allocated + active held <= tower capacity
        Advance time and update traffic
    Settle terminal unused holds; include their waste consistently
    IF training:
        Calculate advantages and return targets
        Update both actors with clipped PPO; fit critic to returns
    ELSE: report service, QoS, switching, waste, and utilization
# Save trained weights after training
\end{lstlisting}

\newpage
\section*{4. Main system formulas, line by line}
Symbols: $C_c$ = tower capacity, $A_c$ = allocated RBs, $R_c$ = active held RBs, $D_u$ = device demand, $M_{u,c}$ = its matching active hold, and $m_u$ = its minimum RB requirement. RB means resource block, the simulator's integer capacity unit.

\textbf{1. Free capacity and safety:}
\begin{equation}
F_c=\max(0,C_c-A_c-R_c),\qquad A_c+R_c\le C_c.
\end{equation}
Active holds reduce free capacity. Pending contracts do not consume RBs.

\textbf{2. Movement prediction:}
\begin{equation}
\widehat{\mathbf x}_u(k)=\mathbf x_u+k\mathbf v_u.
\end{equation}
$\mathbf x_u$ is current position, $\mathbf v_u$ is displacement per interval, and $k$ is the number of forecast intervals. This estimate does not know future random movement.

\textbf{3. Capacity for a new connection:}
\begin{equation}
S_{u,c}=\min(D_u,M_{u,c}+F_c).
\end{equation}
Only a matching active hold can supplement destination free capacity. When evaluating staying, also include the device's own existing allocation at its current tower.

\textbf{4. Admission and QoS:}
\begin{align}
Y_u&=\ind[\text{coverage acceptable}]\ind[S_{u,c}\ge m_u]\ind[\text{radio passes}],\\
Q_u^{\rm HO}&=Y_u\ind[S_{u,c}\ge D_u]\ind[L_{u,c}\le L_u^{\max}].
\end{align}
$\ind[\cdot]$ equals 1 if true, otherwise 0. Admission needs minimum service; handover QoS also needs full demand and acceptable latency. Failed admission does not commit the new association in this design.

\textbf{5. Modeled latency:}
\begin{align}
\rho_c&=(A_c+R_c)/C_c,\qquad s_{u,c}=\max(0,1-S_{u,c}/D_u),\\
L_{u,c}&=L_{\rm base}(1+0.75\rho_c^2+3s_{u,c}).
\end{align}
Load and unmet demand increase estimated latency. This is a synthetic model, not measured packet delay. Base values are 8 ms for eMBB, 2 ms for uRLLC, and 22 ms for mMTC.

\section*{5. Edge cases}
\begin{itemize}
\item \textbf{Strong signal but overloaded:} choose another adequately covered tower, or stay if service is already good.
\item \textbf{Wrong reservation tower:} its hold cannot help elsewhere; release unused active RBs on a handover attempt or expiry and count waste.
\item \textbf{Early arrival:} pending holds cannot help; use free capacity. \textbf{Late or missing arrival:} use only active unexpired holds, otherwise expire them.
\item \textbf{Capacity or radio failure:} keep the old association. If its link has lost coverage, record zero effective service.
\item \textbf{Demand mismatch:} supplement a small hold from free capacity; excess active reserved RBs are waste. Admission can pass while full-demand QoS fails.
\item \textbf{Concurrent arrivals or no alternative:} refresh ledgers after every decision; keep stay legal, but never claim service without capacity and coverage.
\end{itemize}

\newpage
\section*{6. Main reward function}
Measure service to \textbf{all devices}, including those harmed by reservation holds. Let $N$ be device count, $G_u$ its effective grant (zero when uncovered), and $L_u$ the latency estimated from its actual final grant and association.
\begin{align}
F_t&=\frac1N\sum_u\ind[G_u\ge D_u], & &\text{fraction receiving full demand},\\
Q_t&=\frac1N\sum_u\ind[G_u\ge D_u]\ind[L_u\le L_u^{\max}],
& &\text{fraction meeting full-demand QoS},\\
U_t&=\frac{\sum_u\max(0,D_u-G_u)}{\sum_uD_u},
& &\text{fraction of offered demand unmet},\\
O_t&=\frac1N\sum_u\ind[G_u<m_u], & &\text{fraction below minimum service},\\
H_t&=\frac{\text{successful switches in interval}}{N}, & &\text{switching penalty},\\
W_t&=\frac{\text{unused active RBs settled in interval}}{C_{\rm total}},
& &\text{reservation waste penalty},\\
R_t&=\frac{\text{active held RBs before movement}}{C_{\rm total}},
& &\text{holding penalty}.
\end{align}
$C_{\rm total}$ is total capacity across towers. $O_t$ includes both outage and partial under-service. Waste is charged when settled; holding is charged every active interval. Terminal unused holds are included in final reward during training.

\begin{equation}
\boxed{r_t=w_fF_t+w_qQ_t-w_uU_t-w_oO_t-w_hH_t-w_wW_t-w_rR_t.}
\end{equation}
All weights are nonnegative, configurable design parameters. Service and QoS earn reward; unmet demand, under-service, switching, waste, and held resources reduce it. The weights specify the trade-off and are selected before held-out testing. They are not claimed as tuned implementation values.

\section*{7. Learning objective: what PPO optimizes}
\textbf{Overall goal:} learn actor parameters $\theta$ that maximize expected discounted reward:
\begin{equation}
\max_\theta J(\theta),\qquad
J(\theta)=\mathbb E_{\pi_\theta}\left[\sum_{t=0}^{T-1}\gamma^t r_t\right].
\end{equation}
$T$ is episode length and $\gamma$ discounts future rewards. This rewards sustained good connections rather than unnecessary handovers.

\textbf{Advantage:} compare a return target with the critic's prediction:
\begin{equation}
\widehat A_t=\widehat G_t-V_{\rm old}(\mathbf S_t).
\end{equation}
$\widehat G_t$ estimates reward from this interval onward. Positive advantage means better than expected. The training loop computes targets using GAE and normalizes advantages to $A_t$ for the actor update.
\textbf{PPO update:} compare the new and old probabilities of the recorded interval's staged decisions:
\begin{equation}
\omega_t=\prod_{i\in\mathcal D_t}
\frac{\pi_{\theta}(a_{t,i}\mid o_{t,i},m_{t,i})}
{\pi_{\rm old}(a_{t,i}\mid o_{t,i},m_{t,i})}.
\end{equation}
$\mathcal D_t$ includes reservation and connection decisions, each evaluated by its own actor head using stored observations and masks. Compute the product in log space for stability.
\begin{equation}
\max_\theta L_{\rm PPO}=\mathbb E_t\left[
\min\{\omega_t A_t,\clip(\omega_t,1-\epsilon,1+\epsilon)A_t\}\right].
\end{equation}
Clipping limits excessive probability changes; an optional entropy bonus encourages exploration. The critic fits return targets by minimizing $L_V=\mathbb E_t[(\widehat G_t-V_\psi(\mathbf S_t))^2]$. Adam updates the weights. PPO seeks high expected reward, with no guarantee of the best tower on every decision.
\end{document}